\documentclass[11pt]{amsart}
\usepackage[mathscr]{eucal}
%\usepackage[spanish, es-noshorthands]{babel}
\usepackage[english]{babel}
\usepackage{times}
\usepackage{amsmath,amssymb,latexsym,amscd}
\usepackage{booktabs}   
\usepackage{hyperref}
\hypersetup{colorlinks=true,linkcolor=blue,citecolor=brown,linktocpage=true}
%\input xypic
%\xyoption{all}
\usepackage[all,cmtip]{xy}
\usepackage{fancyhdr}
\usepackage{mathalfa}
\usepackage{mathrsfs}
%\usepackage[sans]{dsfont}
\usepackage{upgreek}
%\usepackage{mathpazo}
\usepackage{tikz}
\usepackage{graphicx}
\usepackage{microtype}
\setlength{\emergencystretch}{2em}
\usepackage{tikz-cd}
\usepackage{comment}
\usepackage{todonotes}
% Compatibility fix for todonotes with amsart
\makeatletter
\renewcommand{\listoftodos}[1][\@todonotes@todolistname]{%
  \@starttoc{tdo}{#1}%
}
\makeatother

\newcommand{\done}[1]{%
    \todo[color=green!30]{DONE: #1}%
}
%%%%%%%%%%%%%%%%%%

\DeclareMathOperator{\op}{op}
\DeclareMathOperator{\spec}{spec}
\DeclareMathOperator{\pt}{pt}
\DeclareMathOperator{\tp}{tp}
\DeclareMathOperator{\fd}{fd}
\DeclareMathOperator{\id}{id}
\DeclareMathOperator{\Sp}{sp} 
\DeclareMathOperator{\Ord}{Ord}
\DeclareMathOperator{\Frm}{Frm}
\DeclareMathOperator{\Top}{Top}
\DeclareMathOperator{\Fil}{Fil}
\newcommand{\niton}{\not\owns}

\theoremstyle{plain}
\newtheorem*{theorem*}{Theorem}
\newtheorem{thm}{Theorem}[section]
\newtheorem{cor}[thm]{Corollary}
\newtheorem{lem}[thm]{Lemma}
\newtheorem{prop}[thm]{Proposition}

\theoremstyle{definition}
\newtheorem{dfn}[thm]{Definition}
\newtheorem{obs}[thm]{Observation}
\newtheorem{ej}[thm]{Example}

\renewcommand{\proofname}{Proof}
  %%
    %\providecommand{\corollaryname}{Corollary}
  %\providecommand{\examplename}{Example}
  %\providecommand{\lemmaname}{Lemma}
  %\providecommand{\propositionname}{Proposition}
  %\providecommand{\remarkname}{Remark}
%\providecommand{\theoremname}{Theorem}
%\providecommand{\subexamplename}{Subexample}
  %\providecommand{\definitionname}{Definition}
\begin{document}
%\begin{frontmatter}

\title{Aquí va el título del artículo}


\author{Juan Carlos Monter Cortés and Ángel Zaldívar Corichi}
\email{}
\address{}
 

\address{}
%$\cortext[cor1]{Corresponding author}

\begin{abstract}
Aquí va el resumen
\end{abstract}


\maketitle


\tableofcontents

%Genera una lista automática de todos los TODOs al inicio del documento:

\listoftodos

\section{Introduction}
\todo{Completar la introducción.}
 \done{hay que modificar el formato bastante, hay muchas separaciones en las demostraciones eso se ve raro 
 tambien te suegiero que agrres un paper que te guste y veas su estructura de como escrbiben las cosas}
\cite{picado2021separation, sexton2006ordinal, hochster1969prime}

\cite{simmons2006assembly, sexton2003point, Johnstone82}
\[\begin{tikzcd}
	& {T_2} &&& {\mathbf{(reg)}} & \\
	{\text{S.S.}} & {T_1} & {\text{Neat}} & {\mathbf{(fit)}} & {\mathbf{(sH)}} & {\mathbf{(H)}+\mathbf{(sfit)}} \\
	{\text{S.}} & {T_0} & {\text{Packed}} & {\mathbf{(sfit)}} & {\mathbf{(H)}} & {T_1+\mathbf{(sfit)}} \\
	&& {T_1} && {T_1} & {\text{Spatiality}}
	\arrow[Rightarrow, from=1-2, to=2-1]
	\arrow[Rightarrow, from=1-2, to=2-2]
	\arrow[Rightarrow, from=1-2, to=2-3]
	\arrow[Rightarrow, from=1-5, to=2-4]
	\arrow[Rightarrow, from=1-5, to=2-5]
	\arrow[Rightarrow, from=1-5, to=2-6]
	\arrow[Rightarrow, from=2-1, to=3-1]
	\arrow[Rightarrow, from=2-2, to=3-2]
	\arrow[Rightarrow, from=2-3, to=3-3]
	\arrow[Rightarrow, from=2-4, to=3-4]
	\arrow[Rightarrow, from=2-5, to=3-5]
	\arrow[Rightarrow, from=2-6, to=3-6]
	\arrow[Rightarrow, from=3-3, to=4-3]
	\arrow[Rightarrow, from=3-5, to=4-5]
	\arrow[Rightarrow, from=3-6, to=4-6]
\end{tikzcd}\]

\[\begin{tikzcd}
	& {T_1} & \\
	{\mathbf{(H)}} & KC & {\text{Neat}} \\
	& {\text{S. S.}} & {\text{S.}}
	\arrow[Rightarrow, from=2-1, to=1-2]
	\arrow[Rightarrow, dotted, from=2-1, to=2-2]
	\arrow[Rightarrow, dotted, from=2-1, to=3-2]
	\arrow[Rightarrow, dotted, from=2-2, to=1-2]
	\arrow[Rightarrow, dotted, from=2-2, to=2-3]
	\arrow[Rightarrow, dotted, from=2-2, to=3-2]
	\arrow[Rightarrow, from=2-3, to=1-2]
	\arrow[Rightarrow, dotted, from=2-3, to=3-2]
	\arrow[Rightarrow, dotted, from=2-3, to=3-3]
	\arrow[Rightarrow, from=3-2, to=3-3]
\end{tikzcd}\]

\section{Preliminaries}
\done{Combianar secciones 2 y 3}
\begin{itemize}
	\item Frames, nucleis, $A_j$
	\item $u_a, v_a, w_a$
	\item Points, $\mathcal{U}_*,\mathcal{U}$ spaciality
	\item Open filters, 
	\item Fitted nucleis
	\item HM and HMJ
\end{itemize}

Let $f\colon A\to B$ be a frame morphism with right adjoint $f_*\colon B\to A$. If $F\subseteq A$ and $G\subseteq B$ are filters, we define:
\begin{equation}\label{Imagenfiltros}
b\in f^*F \Longleftrightarrow f_*(b)\in F\qquad a\in f_*G \Longleftrightarrow f(a)\in G,
\end{equation}
where $f^*F$ and $f_*G$ are filters in $B$ and $A$, respectively.

\begin{prop}\label{fF}
For $f=f^*\colon A\to B$ a frame morphism and $G\in B^\wedge$, then $f_*G\in A^\wedge$.
\end{prop}

\begin{proof}
By (\ref{Imagenfiltros}), $f_*G$ is a filter in $A$. We need that $f_*G$ satisfy the open filter condition. Let $X\subseteq A$ be such that $\bigvee X\in f_*G$, with $X$ directed. Then
\[
Y=\{f(x)\mid x\in X\}
\] 
is directed and $f(\bigvee X)=\bigvee f[X]=\bigvee Y\in G$. Since $G$ is an open filter and $f(\bigvee X)=\bigvee f[X]\in G$, there exists $x\in X$ such that $f(x)\in G$. Therefore, $x\in f_*G$.
\end{proof}

\subsection{Compact Quotients and Admissibility Intervals}
\begin{itemize}
	\item $A^\wedge \longleftrightarrow A_j$
	\item $v_F, w_F$ admissibility intervals
	\item 
\end{itemize}

\section{Separation Properties Induced by Compact Quotients}
\subsection{Neat frames}

\begin{dfn}
Neat definition
\end{dfn}

\begin{prop}\label{morfismo} 
Let $(F,Q,G)$ be an HM triple, with 
\[
F\in A^\wedge \quad\text{and}\quad G=\nabla(Q)\in (\mathcal{O}S)^\wedge.
\]
If
\[
k\in[v_G,w_G],
\]
then
\[
\mathcal{U}_*\circ k\circ \mathcal{U}\in[v_F,w_F].
\]
\end{prop}

\begin{proof}
Since $k$ is a nucleus on $\mathcal{O}S$, the composite
\[
\mathcal{U}_*\circ k\circ \mathcal{U}
\]
is a nucleus on $A$. Thus, it is enough to show that
\[
\nabla(\mathcal{U}_*\circ k\circ \mathcal{U})=F.
\]

Let $x\in A$. Since $(F,Q,G)$ is an HM triple,
\[
x\in F
\Longleftrightarrow
Q\subseteq \mathcal{U}(x)
\Longleftrightarrow
\mathcal{U}(x)\in G.
\]

On the other hand, since $k\in[v_G,w_G]$, we have
\[
\nabla(k)=G.
\]
Therefore,
\[
\mathcal{U}(x)\in G
\Longleftrightarrow
k(\mathcal{U}(x))=S.
\]

Using the adjunction $\mathcal{U}\dashv \mathcal{U}_*$ and the fact that $\mathcal{U}(1)=S$, we obtain
\[
k(\mathcal{U}(x))=S
\Longleftrightarrow
\mathcal{U}_*(k(\mathcal{U}(x)))=1.
\]
Hence,
\[
x\in F
\Longleftrightarrow
(\mathcal{U}_*\circ k\circ \mathcal{U})(x)=1
\Longleftrightarrow
x\in\nabla(\mathcal{U}_*\circ k\circ \mathcal{U}).
\]

Thus,
\[
\nabla(\mathcal{U}_*\circ k\circ \mathcal{U})=F,
\]
and consequently
\[
\mathcal{U}_*\circ k\circ \mathcal{U}\in[v_F,w_F].
\]
\end{proof}

Proposition \ref{morfismo} defines a map
\[
\Phi\colon [v_G,w_G]\longrightarrow [v_F,w_F],
\qquad
\Phi(k)=\mathcal{U}_*\circ k\circ \mathcal{U}.
\]
In particular, for every $k\in[v_G,w_G]$, we have
\begin{equation}\label{comparacionvFwF}
v_F\leq \mathcal{U}_*\circ k\circ \mathcal{U}\leq w_F.
\end{equation}

\begin{prop}\label{igualdadwF}
Let $(F,Q,G)$ be an HM triple. Then
\[
w_F=\mathcal{U}_*\circ w_G\circ \mathcal{U}.
\]
\end{prop}

\begin{proof}
By Proposition \ref{morfismo}, since $w_G\in[v_G,w_G]$, we have
\[
\mathcal{U}_*\circ w_G\circ \mathcal{U}\leq w_F.
\]
Thus, it remains to prove the reverse inequality.

By adjunction, it is enough to show that
\[
\mathcal{U}\circ w_F\leq w_G\circ \mathcal{U}.
\]
Let $x\in A$. Since $w_G=[M']$, we have
\[
w_G(\mathcal{U}(x))
=[M'](\mathcal{U}(x))
=(M'\cup \mathcal{U}(x))^\circ.
\]
Hence, we need to prove that
\[
\mathcal{U}(w_F(x))\subseteq (M'\cup \mathcal{U}(x))^\circ.
\]

Since $\mathcal{U}(w_F(x))$ is open, it is enough to show that
\[
\mathcal{U}(w_F(x))\subseteq M'\cup \mathcal{U}(x).
\]
Suppose, by contradiction, that there exists
\[
y\in \mathcal{U}(w_F(x))
\]
such that
\[
y\notin M'
\qquad\text{and}\qquad
y\notin \mathcal{U}(x).
\]
Then $y\in M$ and $x\leq y$. On the other hand,
\[
y\in \mathcal{U}(w_F(x))
\]
implies
\[
w_F(x)\nleq y.
\]
However,
\[
w_F(x)=\bigwedge\{m\in M\mid x\leq m\},
\]
and since $y\in M$ with $x\leq y$, it follows that
\[
w_F(x)\leq y,
\]
a contradiction.

Therefore,
\[
\mathcal{U}(w_F(x))\subseteq (M'\cup \mathcal{U}(x))^\circ,
\]
and hence
\[
\mathcal{U}\circ w_F\leq w_G\circ \mathcal{U}.
\]
By adjunction,
\[
w_F\leq \mathcal{U}_*\circ w_G\circ \mathcal{U}.
\]
Combining both inequalities, we conclude that
\[
w_F=\mathcal{U}_*\circ w_G\circ \mathcal{U}.
\]
\end{proof}

\subsection{Stacked and strongly stacked frames}

\begin{dfn}\label{def:apilado-marcos}
Let $A$ be a frame. We say that $A$ is \emph{stacked} if, for every HM pair $(F,Q)$,
    \[
    (U_*\circ[Q']\circ U)(0)=v_F(0).
    \]
\end{dfn}

\begin{dfn}\label{def:apilado-marcos}
Let $A$ be a frame. We say that $A$ is \emph{strongly stacked} if, for every HM pair $(F,Q)$,
    \[
    U_*\circ[Q']\circ U=v_F;
    \]
\end{dfn}

Clearly, every strongly stacked frame is stacked.
\subsection{KC frames}


\begin{dfn}
Let $A$ be a frame. We say that $A$ is:
\begin{enumerate}
    \item \emph{KC} if every compact quotient of $A$ is closed;
    \item \emph{Hausdorff closed compact}, or \emph{KCH} for short, if every
    compact quotient of $A$ is both closed and Hausdorff.
\end{enumerate}
\end{dfn}

Equivalently, $A$ is KC if, for every $j\in NA$ such that
\[
\nabla(j)\in A^\wedge,
\]
there exists $a\in A$ such that
\[
j=u_a.
\]
Similarly, $A$ is KCH if, in addition, every such compact quotient $A_j$
satisfies property $(H)$.

\section{Comparison with Classical Separation Axioms}
Recall that if a topological space $S$ is $T_1$ and strongly stacked, then
\[
w_G=v_G.
\]
where $G\in \mathcal{O}S^\wedge$. Moreover, every $T_1$ stacked space is sober. This naturally raises the
question of whether analogous results hold in the point-free setting. In this
section, we investigate this question.

\begin{prop}\label{maximosenpt}
Let $A$ be a frame. If
\[
v_F(0)=w_F(0)
\]
for every $F\in A^\wedge$, then every point of $A$ is maximal in $\pt A$.
\end{prop}

\begin{proof}
Let $p,q\in\pt A$ be such that $p\leq q$, and consider the open filter
\[
F_q=\{x\in A\mid x\nleq q\}.
\]
Since $A\setminus F_q=\downarrow q$, it follows directly from the
definition of $w_{F_q}$ that
\begin{equation}\label{igualdadwF_p}
w_{F_q}=w_q.
\end{equation}
In particular,
\[
w_{F_q}(0)=q.
\]

By hypothesis,
\[
q=w_{F_q}(0)=v_{F_q}(0)
=\bigvee_{x\in F_q}\neg x.
\]
For every $x\in F_q$, we have $x\nleq q$. Since $p\leq q$, it follows
that $x\nleq p$. Moreover,
\[
x\wedge\neg x=0\leq p.
\]
Since $p$ is a point, it is meet-irreducible, and therefore
\[
\neg x\leq p.
\]
Thus,
\[
q=v_{F_q}(0)
=\bigvee_{x\in F_q}\neg x
\leq p.
\]
Since $p\leq q$, we conclude that $p=q$. Hence, every point of $A$ is
maximal in 
$\pt A$.
\end{proof}

\begin{prop}\label{focales}
Let $A$ be a frame. The following statements hold:
\begin{enumerate}
    \item If $A$ is $T_1$ and stacked, then
    \[
    v_F(0)=w_F(0)
    \]
    for every $F\in A^\wedge$.

    \item If
    \[
    v_F(0)=w_F(0)
    \]
    for every $F\in A^\wedge$, then $A$ is stacked.

    \item Suppose that
    \[
    v_F(0)=w_F(0)
    \]
    for every $F\in A^\wedge$ and that
    \[
    \{q\in\pt A\mid a\leq q\}\neq\varnothing
    \]
    for every $a\in A$. Then $A$ is $T_1$.
\end{enumerate}
\end{prop}

\begin{proof}
	\mbox{}
\begin{enumerate}
	\item Let $(F,Q)$ be an HM pair. Since $A$ is $T_1$, the space
$S=\pt A$ is $T_1$. Hence $Q=M$. Since $A$ is stacked,
\[
\begin{aligned}
v_F(0)
&=(\mathcal{U}_*\circ[Q']\circ\mathcal{U})(0)\\
&=(\mathcal{U}_*\circ[M']\circ\mathcal{U})(0)\\
&=w_F(0).
\end{aligned}
\]

\item Let $(F,Q)$ be an HM pair. Recall that
\begin{equation}\label{interad}
v_F
\leq
\mathcal{U}_*\circ[Q']\circ\mathcal{U}
\leq
w_F.
\end{equation}
Evaluating at $0$, we obtain
\[
v_F(0)
\leq
(\mathcal{U}_*\circ[Q']\circ\mathcal{U})(0)
\leq
w_F(0).
\]
By hypothesis, $v_F(0)=w_F(0)$, and hence
\[
v_F(0)
=
(\mathcal{U}_*\circ[Q']\circ\mathcal{U})(0).
\]
Therefore, $A$ is stacked.

\item Let $p\in\pt A$ and $a\in A$ be such that
\[
p\leq a.
\]
By hypothesis, there exists $q\in\pt A$ such that
\[
a\leq q.
\]
Hence,
\[
p\leq a\leq q.
\]
By Proposition \ref{maximosenpt}, every point of $A$ is maximal. Therefore, $p=q$, and
consequently
\[
p\leq a\leq p.
\]
Thus, $a=p$, and hence $A$ is $T_1$.
\end{enumerate}
\end{proof}

Thus, in the $T_1$ case, the stacked condition forces the lower and upper
bounds of every admissibility interval to agree at $0$. The strongly stacked
case yields the corresponding equality of the nuclei themselves.

\begin{cor}\label{vF=wF}
Let $A$ be a frame. The following statements hold:
\begin{enumerate}
    \item If $A$ is $T_1$ and strongly stacked, then
    \[
    v_F=w_F
    \]
    for every $F\in A^\wedge$.

    \item If
    \[
    v_F=w_F
    \]
    for every $F\in A^\wedge$, then $A$ is strongly stacked.

    \item Suppose that
    \[
    v_F=w_F
    \]
    for every $F\in A^\wedge$ and that
    \[
    \{q\in\pt A\mid a\leq q\}\neq\varnothing
    \]
    for every $a\in A$. Then $A$ is $T_1$.
\end{enumerate}
\end{cor}

\begin{proof}
\mbox{}
\begin{enumerate}
	\item Let $(F,Q)$ be an HM pair. Since $A$ is strongly stacked,
\[
v_F=\mathcal{U}_*\circ[Q']\circ\mathcal{U}.
\]
Moreover, since $A$ is $T_1$, we have $Q=M$. Hence, by Proposition \ref{igualdadwF},
\[
v_F
=\mathcal{U}_*\circ[M']\circ\mathcal{U}
=w_F.
\]

\item By \eqref{interad},
\[
v_F
\leq
\mathcal{U}_*\circ[Q']\circ\mathcal{U}
\leq
w_F.
\]
Thus, if $v_F=w_F$, then
\[
v_F=\mathcal{U}_*\circ[Q']\circ\mathcal{U},
\]
and therefore $A$ is strongly stacked.

\item Since $v_F=w_F$ implies
\[
v_F(0)=w_F(0),
\]
the result follows from Proposition \ref{focales} (3).
\end{enumerate}
\end{proof}

While the strongly stacked condition yields the collapse of every
admissibility interval, the stacked condition already guarantees this
phenomenon for those associated with completely prime filters. This gives the
following consequence.

\begin{cor}\label{Igualdadnucleos}
Let $A$ be a $T_1$ stacked frame. Then every admissibility interval
associated with a completely prime filter collapses to a singleton. More
precisely, for every $p\in\pt A$,
\[
[v_{F_p},w_{F_p}]=\{u_p\}.
\]
\end{cor}

\begin{proof}
Let $p\in\pt A$. Since $A$ is $T_1$ and stacked, Proposition \ref{focales} gives
\[
v_{F_p}(0)=w_{F_p}(0)=p.
\]
Moreover, as $p$ is maximal,
\[
w_{F_p}=w_p=u_p.
\]

Recall that, for every nucleus $j$, one has
\[
u_{j(0)}\leq j.
\]
Indeed, for every $x\in A$,
\[
j(0)\vee x\leq j(x).
\]
Applying this to $j=v_{F_p}$ yields
\[
u_p=u_{v_{F_p}(0)}\leq v_{F_p}.
\]
On the other hand, by definition of the admissibility interval,
\[
v_{F_p}\leq w_{F_p}=u_p.
\]
Therefore,
\[
[v_{F_p},w_{F_p}]=\{u_p\}.
\]
\end{proof}

\begin{comment}
\textbf{Lo siguiente debe cosiderarse por separado}\\

Para los fines de la investigación que estamos realizando, no puede ocurrir que $S=\emptyset$, donde $S=\pt A$.
Por lo tanto, en nuestro caso siempre suponemos que $S\neq \emptyset$.\\

Notemos que para $a\nleq b$ como antes, se cumple que $\uparrow a$ es un filtro en $A$ y $\downarrow b$ es un ideal. Además, se cumple que 
\[
\uparrow a\, \cap \downarrow b = \emptyset.
\]

Denotamos por $\Fil(A)$ al conjunto de todos los filtros en $A$. Si consideremos el conjunto
\[
\mathcal{F}(a,b)=\{F\in \Fil(A)\mid a\in F\text{ y }b\notin F\},
\]
entonces por Lema de Zorn se cumple que existe un filtro máximo $\mathcal{M}\in \mathcal{F}(a,b)$. Ya que $\mathcal{M}\in \Fil(A)$, podemos construir el núcleo ajustado que admite a $\mathcal{M}$, esto es
\[
\mathcal{M}\subseteq \nabla(v_{\mathcal{M}})
\]
donde $v_\mathcal{M}$ es el núcleo ajustado. Dado que $v_\mathcal{M}\in NA$, tenemos el marco cociente $A_{v_{\mathcal{M}}}$ (y por fines de practicidad lo denotaremos por $A_\mathcal{M}$), cuyo espacio de puntos está dado por
\[
\pt(A_\mathcal{M})=S\setminus \mathcal{M}
\]

Considerando que $S\neq \emptyset$, debemos verificar que ocurre cuando 
\[
\pt(A_\mathcal{M})=\emptyset\quad \text{ y }\quad \pt(A_\mathcal{M})\neq \emptyset.
\]

Si $\pt(A_\mathcal{M})=\emptyset$, entonces $S\subseteq \mathcal{M}$, es decir, para todo $p\in S$, $p \in \mathcal{M}$. De aquí que $p\nleq b$, pues si ocurriera que $p\leq b$, entonces $b\in \mathcal{M}$ y sería una contradicción. \\

Luego, como $\mathcal{M}\in \Fil(A)$, entonces $p\vee b\in\mathcal{M}$ y, al ser $A$ un marco $T_1$, se cumple que
\[
p\vee b=p\quad \text{o}\quad p\vee b=1  \iff b\leq p\quad \text{o}\quad b\nleq p.
\]

Por lo tanto, tenemos los siguientes escenarios:
\[
\text{1) } p\nleq b\;\text{ y }\; b\leq p, \qquad \text{2) } p\nleq b\; \text{ y }\;b\nleq p. 
\]
para todo $p\in S$.\\

Si $\pt(A_\mathcal{M})\neq \emptyset$, existe  $p\in S$ tal que $p\notin \mathcal{M}$.
Luego, como $p\notin \mathcal{M}$, se cumple que $a\nleq p$, pues de lo contrario, si $a\leq p$ eso implicaría que $p\in \mathcal{M}$ y sería una contradicción.\\

Dado que $p\notin \mathcal{M}$, podemos construir el filtro generado por $\langle \mathcal{M}\cup \{p\}\rangle$ y por la maximalidad de $\mathcal{M}$, este filtro satisface que 
\[
b\in \langle\mathcal{M}\cup \{p\}\rangle \quad \text{y}\quad \exists\, m\in \mathcal{M} \text{ tal que }m\wedge p\leq b.
\]

Ya que $v_\mathcal{M}\in NA$, entonces $v_\mathcal{M}(m\wedge p)\leq v_\mathcal{M}(b)$ y por lo tanto $p\leq v_\mathcal{M}(b)$. Al ser $A$ un marco $T_1$, entonces
\[
p=v_\mathcal{M}(b)\quad \text{o}\quad v_\mathcal{M}(b)=1.
\]

De esta manera, si se cumple que $p=v_\mathcal{M}(b)$, entonces $a\nleq p$ y $b\leq p$, es decir, $A$ es espacial. Veamos que $v_\mathcal{M}(b)=1$ nos lleva a una contradicción.\\

\textbf{Hasta aquí llega lo que se debe considerar}\\
\end{comment}

Notice that the additional hypothesis in Proposition \ref{focales} (3) is automatically
satisfied whenever $A$ is spatial. Indeed, if $a<1$ and $A$ is spatial, then
there exists $p\in\pt A$ such that
\[
a\leq p.
\]
Hence,
\[
\mathcal{V}(a)=\{p\in\pt A\mid a\leq p\}\neq\varnothing.
\]
Thus, in the spatial setting, Proposition \ref{IsoVA} reflects the topological fact
that every $T_1$ stacked space is sober.

The point-free situation is, however, genuinely broader, since the $T_1$ and
stacked conditions do not force spatiality. In fact, the following example
shows that this remains true even if stacked is replaced by strongly stacked.

\begin{ej}
Let $B$ be an atomless complete Boolean algebra and consider
\[
A=B\times 2,
\]
where $2=\{0<1\}$. Then $A$ is a complete Boolean algebra. We show that $A$
is $T_1$ and strongly stacked, but not spatial.

First, let
\[
p=(1,0).
\]
We claim that
\[
\pt A=\{p\}.
\]
Indeed, suppose that $(a,b),(c,d)\in A$ satisfy
\[
(a,b)\wedge(c,d)\leq p.
\]
Then
\[
(a\wedge c,b\wedge d)\leq(1,0),
\]
so $b\wedge d=0$. Since $b,d\in 2$, either $b=0$ or $d=0$, and hence
\[
(a,b)\leq p
\qquad\text{or}\qquad
(c,d)\leq p.
\]
Thus, $p\in\pt A$.

Now, let $q=(x,y)\in\pt A$. If $y=0$, then
\[
(1,0)\wedge(x,1)\leq(x,0)=q.
\]
Since $(x,1)\nleq q$, the complete meet-irreducibility of $q$ implies
\[
(1,0)\leq q,
\]
and therefore $x=1$. Hence $q=p$.

If $y=1$, then $x<1$. For $a,b\in B$ such that $a\wedge b\leq x$, we have
\[
(a,1)\wedge(b,1)\leq(x,1)=q.
\]
Since $q$ is a point, either $a\leq x$ or $b\leq x$, showing that
$x\in\pt B$. This is impossible, since an atomless complete Boolean algebra
has no points. Therefore,
\[
\pt A=\{p\}.
\]

Moreover, if
\[
p\leq(a,b),
\]
then $a=1$, and hence $(a,b)$ is either $p$ or $1$. Thus $p$ is maximal and
$A$ is $T_1$.

On the other hand, $A$ is not spatial. Indeed, if $a,b\in B$ with $a\neq b$,
then
\[
(a,0)\neq(b,0),
\]
while
\[
\mathcal{U}(a,0)=\varnothing=\mathcal{U}(b,0).
\]
Hence the spatial reflection
\[
\mathcal{U}\colon A\longrightarrow\mathcal{O}(\pt A)
\]
is not injective.

It remains to prove that $A$ is strongly stacked. Since $B$ is an atomless
complete Boolean algebra, $0$ is its only compact element. Therefore,
\[
K(A)=K(B)\times K(2)=\{(0,0),(0,1)\}.
\]
Since every open filter of a complete Boolean algebra is principal and
$\uparrow a$ is open if and only if $a$ is compact, it follows that
\[
A^\wedge=\{\uparrow(0,0),\uparrow(0,1)\}.
\]
Clearly,
\[
\uparrow(0,0)=A,
\]
whereas
\[
\uparrow(0,1)=B\times\{1\}=F_p.
\]
Hence,
\[
A^\wedge=\{A,F_p\}.
\]

For $F_p$, since $A$ is $T_1$, then
\[
v_{F_p}=w_{F_p}.
\]

For $F=A$, we have
\[
v_F=\operatorname{tp}=w_F.
\]
Consequently,
\[
v_F=w_F
\]
for every $F\in A^\wedge$. By Corollary \ref{vF=wF}, $A$ is strongly stacked.
\end{ej}

\begin{cor}\label{T1apiladonoespacial}
There exist non-spatial $T_1$ strongly stacked frames. In particular, there
exist non-spatial $T_1$ stacked frames.
\end{cor}

The previous example shows that the $T_1$ and stacked conditions do not imply
spatiality. We now compare stackedness with the stronger separation property
$\mathbf{(H)}$. Recall that a frame $A$ satisfies property $\mathbf{(H)}$ if
\[
    1\neq a\nleq b
\]
implies that there exists $c\in A$ such that
\[
    c\nleq a
    \qquad\text{and}\qquad
    \neg c\nleq b.
\]
As shown in \cite{PasekaSmarda1992}, property $\mathbf{(H)}$ implies $T_1$. Nevertheless, property
$\mathbf{(H)}$ does not imply stackedness, as the following example shows.

\begin{ej}
There exists a frame satisfying property $\mathbf{(H)}$ which is not stacked.

Let $I=[0,1]$, endowed with its usual topology, and consider the frame
\[
    A=\mathcal{O}(I).
\]
Since $I$ is a compact Hausdorff space, $A$ is a compact Hausdorff frame.
Moreover, since $I$ is $T_1$, every point of $A$ is maximal and
\[
    D(A)=\pt A
    =
    \{I\setminus\{x\}\mid x\in I\},
\]
where $D(A)$ denotes the set of maximal elements of $A$.

Following the construction of Paseka and \v{S}marda
\cite[Proposition 2.3]{PasekaSmarda1992}, let
\[
    \widetilde{A}
    =
    \{(u,v)\in A\times A_{\neg\neg}\mid u\leq v\},
\]
where
\[
    A_{\neg\neg}=\{a\in A\mid w_0(a)=a\}
\]
is the Boolean algebra of regular elements of $A$. By
\cite[Propositions 2.3 and 2.4]{PasekaSmarda1992}, $\widetilde{A}$ is a compact frame. 
In particular, $\widetilde{A}$ satisfies property $\mathbf{(H)}$.

The maximal elements of $\widetilde{A}$ are precisely the elements of the
form
\[
    D(\widetilde{A})=\{(d,1)\mid d\in D(A)\}.
\]

Since $\widetilde{A}$ satisfies $\mathbf{(H)}$, it is $T_1$, and hence its points
are precisely its maximal elements. Therefore,
\[
    \pt\widetilde{A}
    =
    \{(I\setminus\{x\},I):x\in I\}.
\]

Since $\widetilde{A}$ is compact,
\[
    F=\uparrow 1=\{1\}
\]
is an open filter. Let $(F,Q)$ be the corresponding HM pair. By the
Hofmann--Mislove correspondence,
\[
    Q=\pt\widetilde{A}-F.
\]
Since $1\notin\pt\widetilde{A}$, it follows that
\[
    Q=\pt\widetilde{A}.
\]

Now, since $F=\{1\}$,
\[
    f_F=v_1=\id_{\widetilde{A}},
\]
and therefore
\[
    v_F(0)=0_{\widetilde{A}}.
\]
On the other hand,
\[
\begin{aligned}
    (U_*\circ[Q']\circ U)(0)
    &=
    \bigwedge Q\\
    &=
    \bigwedge\pt\widetilde{A}\\
    &=
    (\varnothing,I)
    \neq
    (\varnothing,\varnothing)
    =
    v_F(0).
\end{aligned}
\]
Thus,
\[
    v_F(0)\neq(U_*\circ[Q']\circ U)(0),
\]
and consequently $\widetilde{A}$ is not stacked.
\end{ej}

Consequently, although property $\mathbf{(H)}$ implies $T_1$, it does not imply
stackedness. Indeed, the preceding example provides a compact Hausdorff
frame which is not stacked.

Together with the strongly stacked non-spatial example above, this
further illustrates the distinction between separation properties and
stackedness in the point-free setting. On the one hand, there exist
non-spatial $T_1$ strongly stacked frames; on the other hand, even the
stronger Hausdorff condition $\mathbf{(H)}$ does not guarantee stackedness.

\section{Quotients and Permanence Properties}

Simmons shows in \cite[Lemma 8.9 and Corollary 8.10]{simmons2004vietoris} that the following diagram
\begin{equation}\label{Qlaxo}
\begin{tikzcd}
	A & A \\
	{\mathcal{O}S} & {\mathcal{O}S}
	\arrow["{f^\infty}", from=1-1, to=1-2]
	\arrow["{U_A}"', from=1-1, to=2-1]
	\arrow["{U_A}", from=1-2, to=2-2]
	\arrow["{F^\infty}"', from=2-1, to=2-2]
\end{tikzcd}
\end{equation}
commutes laxly, that is, $U_A\circ f^\infty\leq F^\infty \circ U_A$. In this diagram $f^\infty$ and $F^\infty$ represent the
fitted nuclei associated with the filters $F\in A^\wedge$ and $\nabla\in \mathcal{O}S^\wedge$, respectively, that is, 
\[
f^\infty=v_F\quad\text{and}\quad F^\infty =v_\nabla.
\]

Here we prove something more general, since we consider $j\in NA$, $F\in A_j^\wedge$ and $j_*F\in A^\wedge$ to produce the diagram
\[\begin{tikzcd}
	A & A \\
	{A_j} & {A_j}
	\arrow["{f^\infty}", from=1-1, to=1-2]
	\arrow["j^*"', from=1-1, to=2-1]
	\arrow["j^*", from=1-2, to=2-2]
	\arrow["{\hat{f}^\infty}"', from=2-1, to=2-2]
\end{tikzcd}\]
where $\hat{f}^\infty$ and $f^\infty$ are the fitted nuclei associated with the filters $j_*F$ and $F$, respectively.

The following lemma is part of the analysis developed in this work and constitutes a key tool in the construction of the commutative diagrams
that follow.

\begin{lem}\label{f1f}
    For $j=j^*$, $f$ and $\hat{f}$ as before, it holds that $j\circ \hat{f}\leq f\circ j$.
\end{lem}
\begin{proof}
    Recall that, by (\ref{Imagenfiltros})
    \[
    \hat{f}=\dot{\bigvee}\{v_y\mid y\in j_*F\} \quad\text{y}\quad f=\dot{\bigvee}\{v_{j(y)}\mid j(y)\in F\}. 
    \]

Let $a\in A$ be. For every $y\in j_*F$, we have
\[
j(v_y(a))=j(y\succ a)\leq (j(y)\succ j(a))=v_{j(y)}(j(a))\leq f(j(a)),
\]
where the last inequality follows from $j(y)\in F$. Therefore, since $j$ preserves arbitrary joins,
\[
j(\hat{f}(a))=j\Big(\bigvee\{v_y(a)\mid y\in j_*F\}\Big)=\bigvee\{j(v_y(a))\mid y\in j_*F\}\leq f(j(a)).
\]

Hence, $j\circ \hat{f}\leq f\circ j$.
\end{proof}

We now prove the above for an arbitrary ordinal.
\begin{cor}\label{finftyf}
    For $j$, $f$ and $\hat{f}$ as before, it holds that $j\circ \hat{f}^\alpha\leq f^\alpha\circ j$ where $\alpha\in \Ord$.
\end{cor}

\begin{proof}
    We proceed by transfinite induction on $\alpha$
    If $\alpha=0$, the result is immediate.

    For the induction step, suppose that the inequality holds for $\alpha$. Then
    \[
    j\circ \hat{f}^{\alpha+1}=j\circ \hat{f}\circ \hat{f}^{\alpha}\leq  f\circ j\circ \hat{f}^\alpha\leq f\circ f^\alpha\circ j=f^{\alpha+1}\circ j,
    \]
    where the first inequality follows from Lemma \ref{f1f} and the second from the induction hypothesis.

    Finally, let $\lambda$ be a limit ordinal. For every $a\in A$,
	\[
	j(\hat{f}(a))=j\Big(\bigvee_{\alpha<\lambda}\hat{f}^\alpha\Big)=\bigvee_{\alpha<\lambda}j(\hat{f}^\alpha(a))\leq \bigvee_{\alpha<\lambda}f^\alpha(j(a))=f^\lambda(j(a)).
	\]
	where the inequality follows from the induction hypothesis. Hence
	\[
	j\circ \hat{f}^\lambda\leq f^\lambda \circ j.
	\]
	This completes the proof.
\end{proof}

\begin{prop}\label{tidyquout}
    If $A$ is a neat frame, then $A_j$ is a neat frame.
\end{prop}

\begin{proof}
Let $A$ be $\alpha$-neat and let $F\in A_j^\wedge$. By Proposition \ref{fF}, $j_*F\in A^\wedge$. Let
\[
\hat{f}=\dot{\bigvee_{y\in j_*F}}v_y\quad\text{and}\quad f=\dot{\bigvee_{x\in F}}v_x.
\]
be the prenuclei associated with $j_*F$ and $F$, respectively.

Since $A$ is neat, there exists an ordinal $\alpha$ such that, for every $y\in j_*F$,
\[
y\vee \hat{d}=1,\quad \text{with }\hat{d}=\hat{f}^\alpha(0).
\]

Let
\[
d=f^\alpha(j(0)),
\]
by Corollary \ref{finftyf},
\[
j(\hat{d})=j(\hat{f}^\alpha(0))\leq f^\alpha(j(0))=d.
\]

Since $j$ is inflationary, $\hat{d}\leq j(\hat{d})$, and hence
\[
\hat{d}\leq d.
\]

Now, let $x\in F$. Thus, $x\in A_j$, we have $j(x)=x$, and therefore $x\in j_*F$. Then,
\[
x\vee \hat{d}=1.
\]

As $\hat{d}\leq d$, it follows that
\[
1=x\vee \hat{d}\leq x\vee d\leq 1.
\]

Hence $x\vee d=1$. Therefore, $A_j$ is $\alpha$-neat.
\end{proof}

The following theorem establishes that tidiness is preserved under binary coproducts.
\begin{prop}\label{tidyplus}
Let $A_1$ and $A_2$ be $\alpha$-neat frames. Then $A_1\oplus A_2$ is $\alpha$-neat.
\end{prop}

\begin{proof}
By construction of the coproduct $A_1\oplus A_2$, consider the canonical frame morphism 
\[
\iota_i\colon A_i\to A_1\oplus A_2, \quad i=1,2
\]
together with their right adjoints $(\iota_i)_*\colon A_1\oplus A_2\to A_i$.

Then, we have the following diagram
\[ 
\begin{tikzcd}
A_i \arrow[r, "\alpha_i"] \arrow[rrr, "\iota_i", bend left] & A_1\times A_2  \arrow[r, "\downarrow(\_)"] & \mathcal{L}(A_1\times A_2) \arrow[r, "\sim"] & A_1\oplus A_2 \arrow[lll, "(\iota_i)_*", bend left]
\end{tikzcd}
\]

By Proposition \ref{fF}, 
\[
(\iota_i)_*F\in A^\wedge_i.
\]

Since $A_i$ is $\alpha$-neat, for every $x_i\in (\iota_i)_*F$ there exists $d_i\in A_i$ such that
\[
x_i\vee d_i=1,
\]
where
\[
d_i=f_i^\alpha(0)\quad \text{and}\quad f_i=\dot{\bigvee_{y\in (\iota_i)_*F}}v_y.
\]

By the definition of $(\iota_i)_*F$,
\[
x_i\in (\iota_i)_*F\iff \iota_i(x_i)\in F.
\]

Thus,
\[
x_1\oplus 1, 1\oplus x_2\in F.
\]

Since $F$ is a filter,
\[
(x_1\oplus 1)\wedge (1\oplus x_2)=x_1\oplus x_2\in F.
\]

Now let
\[
f_F=\dot{\bigvee_{z\in F}}v_z
\]
be the prenucleus associated with $F$. We first show that, for $i=1,2$, 
\[
\iota_i\circ f_i\leq f_F\circ \iota_i.
\]

Indeed, let $a\in A_i$ and $y\in (\iota_i)_*F$. Since $\iota_i(y)\in F$,
\[
\begin{split}
	\iota_i(v_y(a))&=\iota_i(y\succ a)\\
	&\leq (\iota_i(y)\succ \iota_i(a))\\
	&=v_{\iota_i(y)}(\iota_i(a))\\
	&\leq f_F(\iota_i(a)).
\end{split}
\]

Therefore, since the join defining $f_i$ is pointwise and $\iota_i$ preserves arbitrary joins,
\[
\iota_i(f_i(a))=\iota_i\Big(\bigvee_{y\in (\iota_i)_*F}v_y(a)\Big)=\bigvee_{y\in (\iota_i)_*F}\iota_i(v_y(a))\leq f_F(\iota_i(a)).
\]

Hence, 
\[
\iota_i\circ f_i\leq f_F\circ \iota_i.
\]

We now prove by transfinite induction that 
\[
\iota_i\circ f_i^\beta\leq f_F^\beta\circ \iota_i
\]
for every ordinal $\beta$.

If $\beta=0$, the result is immediate. Suppose that it holds for $\beta$. Then
\[
\begin{split}
\iota_i\circ f_i^{\beta +1}&=\iota_i\circ f_i\circ f_i^\beta\\
&\leq f_F\circ \iota_i\circ f_i^\beta\\
&\leq f_F\circ f_F^\beta\circ \iota_i\\
&= f_F^{\beta +1}\circ \iota_i.
\end{split}
\]

Now let $\lambda$ be a limit ordinal. For every $a\in A_i$,
\[
\begin{split}
\iota_i(f_i^\lambda(a))&=\iota_i\Big(\bigvee_{\beta<\lambda}f_i^\beta(a)\Big)\\
&=\bigvee_{\beta<\lambda}\iota_i(f_i^\beta(a))\\
&\leq \bigvee_{\beta<\lambda}f_F^\beta(\iota_i(a))\\
=&f_F^\lambda(\iota_i(a)).
\end{split}
\]

Thus,
\[
\iota_i\circ f_i^\beta\leq f_F^\beta\circ \iota_i
\]
for every ordinal $\beta$.

Evaluating at $0\in A_i$ for $\beta=\alpha$, we obtain
\[
\iota_1(d_1)\leq f_F^\alpha(0\oplus 1)\quad \text{and}\quad \iota_2(d_2)\leq f_F^\alpha(1\oplus 0).
\]

Therefore,
\[
\begin{split}
d_1\oplus d_2&=(d_1\oplus 1)\wedge (1\oplus d_2)\\
&\leq f_F^\alpha(0\oplus 1)\wedge f_F^\alpha(1\oplus 0)\\
&=f_F^\alpha(0\oplus 1\wedge 1\oplus 0).
\end{split}
\]

Set $d=f_F^\alpha(0\oplus 0)$, then $d_1\oplus d_2\leq d$. Finally,
\[
\begin{split}
	1\oplus 1&=(x_1\vee d_1)\oplus (x_2\vee d_2)\\
	&=(x_1\oplus x_2)\vee (d_1\oplus d_2)\\
	\leq (x_1\oplus x_2)\vee d\\
	&\leq 1\oplus 1.
\end{split}
\]

Therefore,
\[
(x_q\oplus x_2)\vee d=1\oplus 1
\]
and $A_1\oplus A_2$ is $\alpha$-tidy.
\end{proof}

\begin{cor}\label{Paso 3}
Let $A_1$ and $A_2$ be $\alpha$-neat and $\beta$-neat, respectively. Then $A_1\oplus A_2$ is $\gamma$-neat, where
\[
\gamma=\max\{\alpha, \beta\}.
\]
\end{cor}

\begin{proof}
Let $\gamma=\max\{\alpha, \beta\}$. Then $A_1$ and $A_2$ are both $\gamma$-neat. Hence, by Proposition \ref{tidyplus}, $A_1\oplus A_2$ is $\gamma$-neat.
\end{proof}

We extend the previous result to arbitrary coproducts.

\begin{prop}\label{coprodE}
Tidiness is preserved under arbitrary coproducts.
\end{prop}

\begin{proof}
Let $\{A_i\}_{i\in I}$ be a family of neat frames, and for each $i\in I$, let $\alpha_i$ be a tidiness degree of $A_i$. Consider 
\[
A=\bigoplus_{i\in I}A_i
\]
and let
\[
\iota_i\colon A_i\to A
\]
denote the canonical coproduct morphisms. Since $\iota_i$ is a frame morphism, it has a right adjoint $(\iota_i)_*$.

Let $F\in A^\wedge$. By Proposition \ref{fF}, 
\[
(\iota_i)_*F\in A_i^\wedge
\]
for every $i\in I$. Moreover,
\[
x_i\in (\iota_i)_*F\iff \iota_i(x_i)\in F.
\]

Set $\alpha=\sup\{\alpha_i\mid i\in I\}$. Then every $A_i$ is $\alpha$-neat. Hence, for every $x_i\in (\iota_i)_*F$,
\[
x_i\vee d_i=1,
\]
where 
\[
d_i=f_i^\alpha(0)\quad \text{and}\quad f_i=\dot{\bigvee_{y\in (\iota_i)_*F}}v_y.
\]

Let
\[
f_F=\dot{\bigvee_x\in F}v_x
\]
be the prenucleus associated with $F$. As in the proof of Proposition \ref{tidyplus}, for every $i\in I$ we have
\[
\iota_i\circ f_i\leq f_F\circ \iota_i,
\]
and, by transfinite induction, 
\[
\iota_i\circ f_i^\beta\leq f_F^\beta\circ \iota_i
\]
for every ordinal $\beta$.

Evaluating at $0\in A_i$ for $\beta=\alpha$, we obtain
\[
\iota_i(d_i)=\iota_i(f^\alpha_i(0))\leq f_F^\alpha(\iota_i(0))=f_F^\alpha(\langle 0\rangle).
\]

Set
\[
d=f_F^\alpha(\langle 0\rangle).
\]

Therefore,
\[
\iota_i(d_i)\leq d
\]
for every $i\in I$. Taking arbitrary joins yields 
\[
\langle d_i\rangle\bigvee_{i\in I}\iota_i(d_i)\leq d.
\]

Now let
\[
\langle x_i\rangle\in F.
\]

Since $x_i\vee d_i=1$ for every $i\in I$, we have
\[
\langle x_i\rangle\vee \langle d_i\rangle=\langle x_i\vee d_i\rangle=\langle 1\rangle.
\]

As $\langle d_i\rangle\leq d$,
\[
\langle 1\rangle=\langle x_i\rangle\vee \langle d_i\rangle\leq \langle x_i\rangle\vee d\leq \langle 1\rangle.
\]

Thus, 
\[
\langle x_i\rangle\vee d=\langle 1\rangle.
\]

Therefore, $A=\bigoplus_{i\in I}A_i$ is $\alpha$-neat, and hence neat.
\end{proof} 

\begin{cor}
The subcategory of neat frames is a coreflective subcategory of $\Frm$.
\end{cor}

Recalling that $v_F=f^\infty$ and $v_{j_*F}=\hat{f}^\infty$, Corollary \ref{finftyf} yields
\[j
\circ v_{j_*F}\leq v_F\circ j
\]

We now compare the compact quotients induced by $F$ and $j_*F$, and construct a frame morphism between them making the corresponding diagram commute.

If we consider $j, v_F$ and $v_{j_*F}$ we have the following diagram
\begin{equation}\label{diagrama H}
\begin{tikzcd}
	A && {A_{j_*F}} \\
	\\
	{A_j} && {A_F}
	\arrow["{(v_{j_*F})^*}", shift left, from=1-1, to=1-3]
	\arrow["j"', from=1-1, to=3-1]
	\arrow["H"', from=1-1, to=3-3]
	\arrow["{(v_{j_*F})_*}", shift left, from=1-3, to=1-1]
	\arrow["h", dotted, from=1-3, to=3-3]
	\arrow["{(v_F)^*}"', shift right, from=3-1, to=3-3]
	\arrow["{(v_F)_*}"', shift right, from=3-3, to=3-1]
\end{tikzcd}
\end{equation}
where $A_F$ and $A_{j_*F}$ are the compact quotients produced by $v_F$ and $v_{j_*F}$, respectively.

Define
\[
H\colon A\to A_F, \qquad H=v_F\circ j.
\]

Since $v_F$ is idempotent, $H(A)\subseteq A_F$. Hence, by restricting $H$ to $A_{j_*F}$, we obtain a map
\[
h\colon A_{j_*F}\to A_F, \qquad h=H|_{A_{j_*F}}.
\]

As finite meets in quotient frames are inherited from the ambient frame, $h$ preserves finite meets. It remains to verify that $h$ preserves arbitrary joins.

\begin{lem}\label{bigvee g}
The map $h\colon A_{j_*F}\to A_F$ preserves arbitrary joins.
\end{lem}

\begin{proof}
	Let $X\subseteq A_{j_*F}$. Since joins in $A_{j_*F}$ are given by
	\[
	\bigvee_{v_{j_*F}}X=v_{j_*F}(\bigvee X)
	\] 
	we have
	\[
	\begin{split}
	h\Big(\bigvee_{v_{j_*F}}X\Big)&=v_F\Big(j\Big(v_{j_*F}\Big(\bigvee X\Big)\Big)\Big)\\
	&\leq v_F\Big(v_F\Big(j\Big(\bigvee X\Big)\Big)\Big)\\
	&=v_F\Big(j\Big(\bigvee X\Big)\Big),
	\end{split}
	\]
	where the inequality follows from Corollary \ref{finftyf} and the equality from the idempotency of $v_F$.

	Since $H$ is a frame morphism, 
	\[
	v_F\Big(j\Big(\bigvee X\Big)\Big)=\bigvee_{v_F}\{H(x)\mid x\in X\}=\bigvee_{v_F}\{h(x)\mid x\in X\}.
	\]
	Therefore,
	\[
	h\Big(\bigvee_{v_{j_*F}}X\Big)\leq \bigvee_{v_F}h[X].
	\]

	Conversely, for every $x\in X$,
	\[
	x\leq \bigvee_{v_{j_*F}}X,
	\]
	and hence, by monotonicity of $h$,
	\[
	h(x)\leq h\Big(\bigvee_{v_{j_*F}}X\Big).
	\]
	Since this holds for every $x\in X$, we have
	\[
	\bigvee_{v_F}h[X]\leq h\Big(\bigvee_{v_{j_*F}}X\Big).
	\]
	Thus,
	\[
	h\Big(\bigvee_{v_{j_*F}}X\Big)= \bigvee_{v_F}h[X].
	\]
	Therefore $h$ preserves arbitrary joins.
\end{proof}

The previous construction yields the following commutative diagram induced by the compact quotients associated with the fitted nuclei.
\begin{prop}\label{VFsquare}
The diagram
\begin{equation}\label{Cuadrado}
    \begin{tikzcd}
	A & {A_{j_*F}} \\
	{A_j} & {A_F}
	\arrow["{v_{j_*F}}", from=1-1, to=1-2]
	\arrow["j"', from=1-1, to=2-1]
	\arrow["h", from=1-2, to=2-2]
	\arrow["{v_F}"', from=2-1, to=2-2]
\end{tikzcd}\end{equation}
is commutative.
\end{prop}

\begin{proof}
	By Lemma \ref{bigvee g}, $h$ is a frame morphism. The commutativity of the diagram follows from the definition of $h$.
\end{proof}

We now study the behavior of strongly stacked frames under quotients.

Recall that, for a frame morphism $f\colon A\to B$, $j\in NA$, and
$k\in NB$, we have
\[
N(f)(j)
=
\bigvee\{u_{f(j(a))}\wedge v_{f(a)}\mid a\in A\}
\in NB
\]
and
\[
N(f)_*(k)=f_*\circ k\circ f\in NA,
\]
where $f_*$ denotes the right adjoint of $f$.

\begin{prop}\label{Cocientefapilado}
Let $A$ be a strongly stacked frame and $j\in NA$. Then $A_j$ is strongly
stacked.
\end{prop}

\begin{proof}
Let $F\in A_j^\wedge$, and set
\[
S=\pt A,
\qquad
S_j=\pt A_j.
\]
Consider the open filter
\[
G=j_*(F)=\{x\in A\mid j(x)\in F\},
\]
and let
\[
Q=S_j\setminus F,
\qquad
R=S\setminus G.
\]
Denote by
\[
\mathcal{U}\colon A\longrightarrow\mathcal{O}S
\qquad\text{and}\qquad
\mathcal{U}_j\colon A_j\longrightarrow\mathcal{O}S_j
\]
the corresponding spatial reflections.

Set
\[
l=N(\mathcal{U}_j)_*([Q'])
=(\mathcal{U}_j)_*\circ[Q']\circ\mathcal{U}_j
\]
and
\[
m=N(j^*)(v_G).
\]
By Proposition \ref{morfismo},
\[
v_F\leq l.
\]
Thus, it is enough to prove
\[
l\leq m\leq v_F.
\]

We first prove that $m\leq v_F$. Since $G=j_*(F)$, the fitted nuclei
associated with $G$ and $F$ satisfy
\begin{equation}\label{eq1}
j^*\circ v_G\leq v_F\circ j^*.
\end{equation}
Moreover,
\[
m=N(j^*)(v_G)
=
\bigvee_{x\in A}
\left(
u_{j^*(v_G(x))}\wedge v_{j^*(x)}
\right).
\]
By \eqref{eq1}, for every $x\in A$,
\[
j^*(v_G(x))\leq v_F(j^*(x)),
\]
and hence
\[
u_{j^*(v_G(x))}
\leq
u_{v_F(j^*(x))}.
\]
Therefore,
\[
u_{j^*(v_G(x))}\wedge v_{j^*(x)}
\leq
u_{v_F(j^*(x))}\wedge v_{j^*(x)}
\leq v_F.
\]
Taking joins over $x\in A$, we obtain
\begin{equation}\label{eq2}
m\leq v_F.
\end{equation}

We now prove that $l\leq m$. Since $A$ is strongly stacked,
\[
v_G=N(\mathcal{U})_*([R'])
=\mathcal{U}_*\circ[R']\circ\mathcal{U}.
\]
Hence,
\[
m
=
N(j^*)\bigl(N(\mathcal{U})_*([R'])\bigr).
\]

To compare the $u$-components in the decompositions of $l$ and $m$, consider
the auxiliary nucleus
\[
n=N(\mathcal{U}_j\circ j^*)_*([Q']).
\]
We first observe that, for every $x\in A_j$,
\[
j^*(n(x))=l(x).
\]
Indeed,
\[
\begin{aligned}
j^*(n(x))
&=
j^*\left(
j_*\circ(\mathcal{U}_j)_*
\circ[Q']\circ\mathcal{U}_j\circ j^*
\right)(x)\\
&=
j^*j_*\left(
\bigvee
\left\{
a\in A_j
\;\middle|\;
\mathcal{U}_j(a)
\subseteq
Q'\cup\mathcal{U}_j(j^*(x))
\right\}
\right)\\
&=
\bigvee
\left\{
a\in A_j
\;\middle|\;
\mathcal{U}_j(a)
\subseteq
Q'\cup\mathcal{U}_j(x)
\right\}\\
&=l(x).
\end{aligned}
\]

Now, under the usual identification of $S_j$ with the corresponding subspace
of $S$, we have
\[
Q=R\cap S_j.
\]
Using the decomposition of a nucleus in terms of its $u$- and $v$-components,
we have
\[
l
=
\bigvee_{x\in A_j}
\left(
u_{l(x)}\wedge v_x
\right)
=
\bigvee_{x\in A_j}
\left(
u_{j^*(n(x))}\wedge v_x
\right).
\]

On the other hand, since
\[
m
=
N(j^*)\bigl(N(\mathcal{U})_*([R'])\bigr),
\]
we have
\[
m
=
\bigvee_{a\in A}
\left(
u_{j^*(N(\mathcal{U})_*([R'])(a))}
\wedge
v_{j^*(a)}
\right).
\]

For $x\in A_j$, the explicit description of the right adjoint gives
\[
\begin{aligned}
j^*\bigl(N(\mathcal{U})_*([R'])(x)\bigr)
&=
j^*\left(
\bigvee
\left\{
a\in A
\;\middle|\;
\mathcal{U}(a)
\subseteq
R'\cup\mathcal{U}(x)
\right\}
\right)\\
&=
\bigvee
\left\{
j^*(a)\in A_j
\;\middle|\;
\mathcal{U}(j^*(a))
\subseteq
Q'\cup\mathcal{U}(x)
\right\}.
\end{aligned}
\]
Moreover, since $x\in A_j$,
\[
\mathcal{U}(j^*(a))
\subseteq
Q'\cup\mathcal{U}(x)
\]
is equivalent, after restriction to $S_j$, to
\[
\mathcal{U}_j(j^*(a))
\subseteq
Q'\cup\mathcal{U}_j(x).
\]
Consequently,
\[
j^*(n(x))
\leq
j^*\bigl(N(\mathcal{U})_*([R'])(x)\bigr).
\]
Therefore,
\[
u_{j^*(n(x))}
\leq
u_{j^*(N(\mathcal{U})_*([R'])(x))}
\]
for every $x\in A_j$, and hence
\[
u_{l(x)}\wedge v_x
=
u_{j^*(n(x))}\wedge v_x
\leq
u_{j^*(N(\mathcal{U})_*([R'])(x))}\wedge v_x
\leq m.
\]
Taking joins over $x\in A_j$, we obtain
\[
l\leq m.
\]

Combining this inequality with \eqref{eq2} and $v_F\leq l$, we obtain
\[
v_F\leq l\leq m\leq v_F.
\]
Hence,
\[
v_F=l
=
(\mathcal{U}_j)_*\circ[Q']\circ\mathcal{U}_j.
\]
Therefore, $A_j$ is strongly stacked.
\end{proof}

Since every strongly stacked frame is stacked, we obtain the following result.

\begin{cor}
Let $A$ be a strongly stacked frame and $j\in NA$. Then $A_j$ is stacked.
\end{cor}

\begin{proof}
By Proposition \ref{Cocientefapilado}, $A_j$ is strongly stacked, and hence stacked.
\end{proof}

We first observe that the KC property is inherited by quotients.

\begin{prop}\label{KCquout}
Let $A$ be a KC frame and $j\in NA$. Then $A_j$ is KC.
\end{prop}

\begin{proof}
Let $k\in N(A_j)$ be such that $(A_j)_k$ is compact. Then
\[
\nabla(k)\in A_j^\wedge.
\]
By Proposition \ref{fF},
\[
j_*(\nabla(k))\in A^\wedge.
\]
Consider the nucleus
\[
l=j_*\circ k\circ j^*\in NA.
\]
Since
\[
\nabla(l)=j_*(\nabla(k)),
\]
the quotient $A_l$ is compact. Since $A$ is KC, there exists $a\in A$ such that
\[
l=u_a.
\]

We claim that
\[
k=u_{j^*(a)}.
\]
Indeed, from $l=u_a$ we have
\[
a\vee x
=
l(x)
=
j_*\bigl(k(j^*(x))\bigr)
\]
for every $x\in A$. Applying $j^*$ to both sides yields
\[
j^*(a)\vee j^*(x)
=
k(j^*(x)),
\]
since $j^*\circ j_*=\id_{A_j}$.

Now let $x\in A_j$. Since $j^*\colon A\to A_j$ is surjective, we have
$j^*(x)=x$, and therefore
\[
\begin{aligned}
u_{j^*(a)}(x)
&=j^*(a)\vee x\\
&=j^*(a)\vee j^*(x)\\
&=k(j^*(x))\\
&=k(x).
\end{aligned}
\]
Hence
\[
k=u_{j^*(a)}.
\]
Thus $(A_j)_k$ is closed. Therefore, $A_j$ is KC.
\end{proof}

\begin{prop}\label{KCT1}
Every KC frame is $T_1$.
\end{prop}

\begin{proof}
Let $p\in\pt A$. The filter
\[
F_p=\nabla(w_p)=\{x\in A\mid x\nleq p\}
\]
is completely prime and, in particular, $F_p\in A^\wedge$. Hence $A_{w_p}$
is a compact quotient of $A$. Since $A$ is KC, this quotient is closed.
Therefore,
\[
w_p=u_a
\]
for some $a\in A$. Evaluating at $0$, we obtain
\[
a=u_a(0)=w_p(0)=p,
\]
and hence
\[
w_p=u_p.
\]

Now let $a\in A$ be such that
\[
p<a\leq1.
\]
Since $a\nleq p$, we have
\[
w_p(a)=1.
\]
On the other hand,
\[
u_p(a)=p\vee a=a.
\]
Since $w_p=u_p$, it follows that
\[
a=1.
\]
Thus $p$ is maximal. Since every point of $A$ is maximal, $A$ is $T_1$.
\end{proof}


\section{The Spatial Case}

\begin{lem}\label{inducciontransfinita}
Let $\alpha\in \Ord$ and $X\in CS$. Then
\begin{equation}\label{eq:induccion}
\partial_F^\alpha(X)=\big(f_F^\alpha(X')\big)'.
\end{equation}
\end{lem}

\begin{proof}
We proceed by transfinite induction on \(\alpha\).

For \(\alpha=0\), both \(\partial_F^0\) and \(f_F^0\) are the identity. Hence,

$$
\partial_F^0(X)
=X
=(X')'
=\left(f_F^0(X')\right)'.
$$

Suppose that the result holds for some ordinal \(\alpha\). Then

$$
\begin{aligned}
\partial_F^{\alpha+1}(X)
&=\partial_F\left(\partial_F^\alpha(X)\right)\\
&=\left(f_F\left(\left(\partial_F^\alpha(X)\right)'\right)\right)'\\
&=\left(f_F\left(f_F^\alpha(X')\right)\right)'\\
&=\left(f_F^{\alpha+1}(X')\right)'.
\end{aligned}
$$

Finally, let \(\lambda\) be a limit ordinal and suppose that the result holds for every \(\alpha<\lambda\). By definition and De Morgan's law,

$$
\begin{aligned}
\left(\partial_F^\lambda(X)\right)'
&=
\left(\bigcap_{\alpha<\lambda}\partial_F^\alpha(X)\right)'\\
&=
\bigcup_{\alpha<\lambda}
\left(\partial_F^\alpha(X)\right)'\\
&=
\bigcup_{\alpha<\lambda}
f_F^\alpha(X')\\
&=
f_F^\lambda(X'),
\end{aligned}
$$

where the third equality follows from the induction hypothesis. Therefore,

$$
\partial_F^\lambda(X)
=
\left(f_F^\lambda(X')\right)'.
$$

This completes the proof.
\end{proof}

The following is an immediate consequence of Lemma \ref{inducciontransfinita}.


\begin{cor}\label{cor:infinito}
Let $U\in\mathcal{O}S$. Then
\begin{equation}\label{eq:infinito}
    \begin{split}
\left(\partial_F^\infty(U')\right)'=v_F(U).
	\end{split}
\end{equation}
\end{cor}


A particular instance of diagram (\ref{Cuadrado}) is obtained by taking $j^*=\mathcal{U}_A$, where $S=\pt A$. Thus, we obtain 
the following commutative diagram:
\[\begin{tikzcd}
	A & {A_F} \\
	{\mathcal{O}S} & {\mathcal{O}S_\nabla}
	\arrow["{v_F}", from=1-1, to=1-2]
	\arrow["{\mathcal{U}_A}"', from=1-1, to=2-1]
	\arrow["h", from=1-2, to=2-2]
	\arrow["{v_\nabla}"', from=2-1, to=2-2]
\end{tikzcd}\]

Recall that $\pt A_F=Q$. Hence, the spatial reflection of $A_F$ is given by
\[
\mathcal{U}_{A_F}\colon A_F\to \mathcal{O}\pt A_F=\mathcal{O}Q.
\]
 
Similarly, since $\pt((\mathcal{O}S)_\nabla)=Q$, we have the spatial reflection
\[
\mathcal{U}_{\mathcal{O}S_\nabla}\colon \mathcal{O}S_\nabla\to \mathcal{O}\pt(\mathcal{O}S)_\nabla=\mathcal{O}Q.
\]

Consequently, the two induced morphisms from $A$ to $\mathcal{O}Q$ are both given by
\[
x\mapsto \mathcal{U}_A(x)\cap Q.
\]

Therefore, the preceding constructions fit into the following commutative diagram:
\begin{equation}\label{Qcuadrado}
\begin{tikzcd}
	A && {A_F} & \\
	&&& {\mathcal{O}Q} \\
	{\mathcal{O}S} && {\mathcal{O}S_\nabla}
	\arrow["{v_F}", from=1-1, to=1-3]
	\arrow["{\mathcal{U}_A}"', from=1-1, to=3-1]
	\arrow["{\mathcal{U}_{A_F}}", from=1-3, to=2-4]
	\arrow["h"', from=1-3, to=3-3]
	\arrow["{v_\nabla}"', from=3-1, to=3-3]
	\arrow["{\mathcal{U}_{\mathcal{O}S_\nabla}}"', from=3-3, to=2-4]
\end{tikzcd}\end{equation}

Since the above diagram is determined by $F\in A^\wedge$ and the associated compact saturated subset $Q\in \mathcal{Q}S$, 
we refer to diagram (\ref{Qcuadrado}) as the \emph{$Q$-square}.

\begin{thm}\label{haus1}
Under the hypotheses of the $Q$-square, if $A$ satisfies property $\mathbf{(H)}$, then
\[
\mathcal{O}Q\cong \uparrow Q'.
\]
Equivalently, the frame of open sets of the point space of $A_F$ is isomorphic to a compact closed quotient of 
a Hausdorff space.
\end{thm}

\begin{proof}
Suppose that $A$ satisfies $\mathbf{(H)}$. Since $\mathbf{(H)}$ is conservative, $S=\pt A$ is a $T_2$ space. Hence, for every $Q\in \mathcal{Q}S$, the set $Q$ is closed, and therefore $Q'\in \mathcal{O}S$.

Moreover, since every $T_2$ space is strongly stacked, we have
\[
v_\nabla=[Q'].
\]
As $Q'\in \mathcal{O}S$, it follows that 
\[
[Q']=u_{Q'}
\]

Consequently,
\[
(\mathcal{O}S)_{v_\nabla}=(\mathcal{O}S)_{u_{Q'}}\cong \uparrow Q'.
\]

Now, $u_{Q'}$ is a complemented nucleus in $\mathcal{O}S$, by \cite[Proposition VI.3.3]{PicadoPultr}, implies that $\mathcal{O}S_{v_\nabla}$ is spatial. Therefore, 
\[
(\mathcal{O}S)_{v_\nabla}\cong\mathcal{O}\pt((\mathcal{O}S)_{v_\nabla}).
\]
Since
\[
\pt((\mathcal{O}S)_{v_\nabla})=Q,
\]
we obtain
\[
\mathcal{O}Q\cong (\mathcal{O}S)_{v_\nabla}\cong \uparrow Q'.
\]
\end{proof}

\begin{lem}\label{lem:auxiliar}
Let $A=\mathcal{O}S$ be spatial and let $E\subseteq S$. Then
\[
\mathcal{U}_*\circ [E]\circ \mathcal{U}=[E].
\]
\end{lem}

\begin{proof}
Let $V\in\mathcal{O}S$. Since
\[
\mathcal{U}\colon \mathcal{O}S\longrightarrow
\mathcal{O}\pt(\mathcal{O}S)
\]
is an isomorphism, there exist $x,y\in\mathcal{O}S$ such that
\[
V=\mathcal{U}(y)
\]
and
\[
(E\cup V)^\circ=\mathcal{U}(x).
\]

Then
\[
\begin{aligned}
(\mathcal{U}_*\circ [E]\circ \mathcal{U})(y)
&=\mathcal{U}_*\bigl([E](\mathcal{U}(y))\bigr)\\
&=\mathcal{U}_*\bigl((E\cup V)^\circ\bigr)\\
&=\mathcal{U}_*(\mathcal{U}(x))\\
&=x.
\end{aligned}
\]

On the other hand,
\[
\mathcal{U}([E](y))
=(E\cup \mathcal{U}(y))^\circ
=(E\cup V)^\circ
=\mathcal{U}(x).
\]
Since $\mathcal{U}$ is injective, it follows that
\[
[E](y)=x.
\]

Therefore,
\[
(\mathcal{U}_*\circ [E]\circ \mathcal{U})(y)
=[E](y)
\]
for every $y\in\mathcal{O}S$. Hence,
\[
\mathcal{U}_*\circ [E]\circ \mathcal{U}=[E].
\]
\end{proof}

We now show that both notions are conservative; namely, for every topological space $S$,
\[
S \text{ is stacked}
\Longleftrightarrow
\mathcal{O}S \text{ is stacked},
\]
and
\[
S \text{ is strongly stacked}
\Longleftrightarrow
\mathcal{O}S \text{ is strongly stacked}.
\]

\begin{prop}\label{apiladoconservativa}
The notions of strongly stacked and stacked are conservative.
\end{prop}

\begin{proof}
Let $S\in\mathbf{Top}$ and set $A=\mathcal{O}S$.

First, suppose that $S$ is strongly stacked. Then, for every HM pair
$(F,Q)$,
\[
v_F=[Q'].
\]
By Lemma~5.3,
\[
\mathcal{U}_*\circ[Q']\circ\mathcal{U}=[Q'].
\]
Therefore,
\[
\mathcal{U}_*\circ[Q']\circ\mathcal{U}=v_F,
\]
and hence $\mathcal{O}S$ is strongly stacked.

Now, suppose that $S$ is stacked. Then, for every HM pair $(F,Q)$,
\[
Q=\partial_F^\infty(S).
\]
Hence,
\[
Q'=\left(\partial_F^\infty(S)\right)'.
\]
By Corollary~5.2,
\[
v_F(\varnothing)
=\left(\partial_F^\infty(S)\right)'
=Q'.
\]
In particular, $Q'$ is open. By Lemma~5.3,
\[
\mathcal{U}_*\circ[Q']\circ\mathcal{U}=[Q'].
\]
Therefore,
\[
\begin{aligned}
(\mathcal{U}_*\circ[Q']\circ\mathcal{U})(\varnothing)
&=[Q'](\varnothing)\\
&=(Q')^\circ\\
&=Q'\\
&=v_F(\varnothing).
\end{aligned}
\]
Thus, $\mathcal{O}S$ is stacked.

Since all the implications used above are equivalences, the converse statements follow by reversing the argument.
\end{proof}

Recall that $A^\wedge$ is endowed with the topology generated by the basic
open sets
\[
D(x)=\{F\in A^\wedge\mid x\in F\},
\qquad x\in A.
\]
Moreover, let $VA$ denote the Vietoris modification of $A$ introduced in
\cite{simmons2004vietoris}. The points of $VA$ are of the form
\[
(F,a),
\qquad
F\in A^\wedge,\quad
a\in[v_F(0),w_F(0)].
\]
Consider the continuous surjection
\[
\pi\colon\pt(VA)\longrightarrow A^\wedge,
\qquad
(F,a)\longmapsto F.
\]
It induces the frame morphism
\[
\mathcal{O}(\pi)=\pi^{-1}\colon
\mathcal{O}(A^\wedge)\longrightarrow
\mathcal{O}(\pt(VA)).
\]

Recall also that the subbasic open sets of $\pt(VA)$ are determined by
\begin{equation}\label{fibra}
(F,a)\in\Box x
\Longleftrightarrow
x\in F,
\qquad
(F,a)\in\Diamond x
\Longleftrightarrow
x\nleq a.
\end{equation}

Particularly,
\begin{equation}\label{fibra2}
a\in [F](x)\iff x\in F\quad \text{and}\quad a\in \langle F\rangle(x)\iff x\nleq a.
\end{equation}

\begin{prop}\label{IsoVA}
Let $A$ be a frame, and consider the continuous surjection
\[
\pi\colon\pt(VA)\longrightarrow A^\wedge,
\qquad
(F,a)\longmapsto F.
\]
Then the following statements hold:
\begin{enumerate}
    \item If $A$ is $T_1$ and stacked, then
    \[
    \mathcal{O}(\pi)=\pi^{-1}\colon
    \mathcal{O}(A^\wedge)\longrightarrow\mathcal{O}(\pt(VA))
    \]
    is an isomorphism.

    \item If $\mathcal{O}(\pi)$ is an isomorphism, then $A$ is stacked.

    \item If $\mathcal{O}(\pi)$ is an isomorphism and
    \[
    \mathcal{V}(a)=\{p\in\pt A\mid a\leq p\}\neq\varnothing
    \]
    for every $a<1$, then $A$ is $T_1$.
\end{enumerate}
\end{prop}

\begin{proof}
	\mbox{}
	\begin{enumerate}
		\item Suppose that $A$ is $T_1$ and stacked. By Proposition \ref{focales},
\[
v_F(0)=w_F(0)
\]
for every $F\in A^\wedge$. Hence, each fiber of
\[
\pi\colon\pt(VA)\longrightarrow A^\wedge
\]
contains a unique element. Therefore, $\pi$ is bijective.

It remains to prove that
\[
\pi^{-1}\colon A^\wedge\longrightarrow\pt(VA)
\]
is continuous. Recall that the subbasic open sets of $\pt(VA)$ are given by
\[
(F,a)\in[-](x)\Longleftrightarrow x\in F
\]
and
\[
(F,a)\in\langle-\rangle(x)\Longleftrightarrow x\nleq a.
\]

For the first family,
\[
(F,a)\in[-](x)
\Longleftrightarrow
x\in F
\Longleftrightarrow
F\in D(x),
\]
and therefore
\[
\pi([-](x))=D(x).
\]

Now consider $\langle-\rangle(x)$. Since the fibers of $\pi$ collapse, for every
$(F,a)\in\pt(VA)$ we have
\[
a=v_F(0).
\]
Thus,
\[
\begin{aligned}
(F,a)\notin\langle-\rangle(x)
&\Longleftrightarrow x\leq a\\
&\Longleftrightarrow x\leq v_F(0)\\
&\Longleftrightarrow v_F(x)=v_F(0)\\
&\Longleftrightarrow x\notin F.
\end{aligned}
\]
Hence,
\[
(F,a)\in\langle-\rangle(x)
\Longleftrightarrow
x\in F
\Longleftrightarrow
F\in D(x),
\]
and therefore
\[
\pi(\langle-\rangle(x))=D(x).
\]

Thus, the image under $\pi$ of every subbasic open set of $\pt(VA)$ is open in
$A^\wedge$. Hence $\pi^{-1}$ is continuous, and consequently $\pi$ is a
homeomorphism. Therefore,
\[
\mathcal{O}(\pi)=\pi^{-1}\colon
\mathcal{O}(A^\wedge)\longrightarrow\mathcal{O}(\pt(VA))
\]
is an isomorphism.

\item Suppose that
\[
\mathcal{O}(\pi)=\pi^{-1}\colon
\mathcal{O}(A^\wedge)\longrightarrow\mathcal{O}(\pt(VA))
\]
is an isomorphism. Since $\pi^{-1}$ is bijective, $\pi$ is injective. For every
$F\in A^\wedge$,
\[
(F,v_F(0)),(F,w_F(0))\in\pt(VA)
\]
and
\[
\pi(F,v_F(0))=F=\pi(F,w_F(0)).
\]
Hence,
\[
v_F(0)=w_F(0).
\]
By Proposition \ref{focales}, $A$ is stacked.

\item Suppose that $\mathcal{O}(\pi)$ is an isomorphism and that
\[
\mathcal{V}(a)=\{q\in\pt A\mid a\leq q\}\neq\varnothing
\]
for every $a<1$. By part~(2),
\[
v_F(0)=w_F(0)
\]
for every $F\in A^\wedge$. Therefore, by Proposition \ref{maximosenpt}, every point of $A$ is
maximal.

Let $p\in\pt A$ and $a\in A$ be such that
\[
p\leq a.
\]
By hypothesis, there exists $q\in\pt A$ such that
\[
a\leq q.
\]
Thus,
\[
p\leq a\leq q.
\]
Since $p$ and $q$ are maximal, we have $p=q$, and therefore
\[
p\leq a\leq p.
\]
Hence $a=p$, and $A$ is $T_1$.
\end{enumerate}
\end{proof}


\section{Relations among Compact-Quotient Separation Properties}
\todo{Ejemplo de neat no KC}
\todo{Comprobar si neat implica stacked}
\todo{KC implica S.S.?}
\todo{Existe alguna propiedad P tal que (H)+P implica staked}
\todo{T1+S.+P implica (H)?}
\todo{Relacionar S. y S.S. con fit y sfit}
\todo{Relacionar el trabajo con lo presentado por Guram y Sebastian}

For every HM triple $(F,Q,G)$, we have
\[
\mathcal{U}\circ v_F\leq v_G\circ\mathcal{U}.
\]
Suppose now that $A$ is strongly stacked. Then
\[
v_F=\mathcal{U}_*\circ[Q']\circ\mathcal{U}.
\]
Since
\[
v_G\leq[Q'],
\]
we obtain
\[
\mathcal{U}_*\circ v_G\circ\mathcal{U}
\leq
\mathcal{U}_*\circ[Q']\circ\mathcal{U}
=v_F.
\]
Together with the previous inequality, this yields
\[
v_F=\mathcal{U}_*\circ v_G\circ\mathcal{U},
\]
or equivalently,
\[
\mathcal{U}\circ v_F
=
v_G\circ\mathcal{U}.
\]
Thus, if $A$ is strongly stacked, the corresponding diagram is commutative.

\begin{prop}\label{EficienteFapilado}
Let $A$ be a neat frame. If
\[
\mathcal{U}\circ v_F
=
v_\nabla\circ\mathcal{U}
\]
for every HM triple $(F,Q,\nabla)$, then $A$ is strongly stacked.
\end{prop}

\begin{proof}
Let $S=\pt A$ and let $(F,Q,\nabla)$ be an HM triple.

Since $A$ is neat, $\mathcal{O}S$ is neat. Hence $S$ is stacked and packed.
In particular, every compact saturated subset $Q\subseteq S$ is closed, so
\[
Q'\in\mathcal{O}S.
\]

Since $S$ is stacked,
\[
Q=\partial_\nabla^\infty(S),
\]
and therefore, by Corollary \ref{cor:infinito},
\[
v_\nabla(\varnothing)=Q'.
\]
Because $\mathcal{O}S$ is neat, it follows that
\[
v_\nabla=u_{Q'}=[Q'].
\]

By hypothesis,
\[
\mathcal{U}\circ v_F
=
v_\nabla\circ\mathcal{U}
=
[Q']\circ\mathcal{U}.
\]
Applying the right adjoint $\mathcal{U}_*$, we obtain
\[
v_F
=
\mathcal{U}_*\circ[Q']\circ\mathcal{U}.
\]
Therefore, $A$ is strongly stacked.
\end{proof}

Consequently, we have the following result.

\begin{cor}
If $A$ is a neat frame, then $\pt A$ is strongly stacked.
\end{cor}

\begin{proof}
Let $S=\pt A$. Since $A$ is neat, $\mathcal{O}S$ is tidy, and therefore
$S$ is strongly stacked.
\end{proof}

\begin{prop}
Let $A$ be a KC-frame. Then $A$ is neat.
\end{prop}

\begin{proof}
Let $F\in A^\wedge$ and put
\[
d=v_F(0).
\]
Since $A$ is KC and
\[
v_F\in[v_F,w_F],
\]
there exists $a\in A$ such that
\[
v_F=u_a.
\]
Evaluating at $0$, we obtain
\[
d=v_F(0)=u_a(0)=a.
\]
Hence
\[
v_F=u_d.
\]
Since this holds for every $F\in A^\wedge$, the frame $A$ is
neat.
\end{proof}

\begin{prop}
Let $A$ be a KC-frame. Then $A$ is stacked if and only if it is
strongly stacked.
\end{prop}

\begin{proof}
Let $(F,Q)$ be an HM-pair and put
\[
d=v_F(0)
\qquad\text{and}\qquad
a=w_F(0).
\]
Since every KC-frame is $T_1$, we have
\[
\mathcal{U}_*\circ[Q']\circ\mathcal{U}=w_F.
\]
Moreover, by the previous proposition,
\[
v_F=u_d.
\]
Since $A$ is KC, $w_F$ is closed as well, and therefore
\[
w_F=u_a.
\]

Consequently,
\[
\begin{aligned}
A\text{ is strongly stacked}
&\iff v_F=w_F\\
&\iff u_d=u_a\\
&\iff d=a\\
&\iff v_F(0)=w_F(0)\\
&\iff A\text{ is stacked}.
\end{aligned}
\]
Since this holds for every HM-pair $(F,Q)$, the result follows.
\end{proof}

\todo{Agregar como corolario el caso espacial del resultado anterior.}

\begin{cor}
Let $A$ be a spatial frame. Then $A$ is KC if and only if it is
neat.
\end{cor}

\begin{proof}
Suppose first that $A$ is KC. By the previous results, every KC-frame
is neat.

Conversely, let $A$ be neat. Since $A$ is spatial, we may write
\[
A\cong\mathcal{O}(S)
\]
for some topological space $S$. By the spatial characterization of
neatnees, $S$ is packed and stacked. Moreover, neatness implies
that $S$ is strongly stacked. Hence $S$ is packed and strongly
stacked, and therefore $\mathcal{O}(S)$ is KC. Thus $A$ is KC.
\end{proof}

\begin{prop}
Let $A$ be a KC-frame. If $A$ is subfit, then $A$ is strongly
stacked.
\end{prop}

\begin{proof}
Let $(F,Q)$ be an HM-pair and put
\[
d=v_F(0)
\qquad\text{and}\qquad
a=w_F(0).
\]
Since $A$ is KC, both $v_F$ and $w_F$ are closed nuclei. Hence
\[
v_F=u_d
\qquad\text{and}\qquad
w_F=u_a.
\]
By subfitness and \cite[Theorem~3.5]{simmons2006regularity}, we have
\[
d=a.
\]
Therefore,
\[
v_F=u_d=u_a=w_F.
\]
Since $A$ is $T_1$, we also have
\[
w_F=\mathcal{U}_*\circ[Q']\circ\mathcal{U}.
\]
Consequently,
\[
v_F=\mathcal{U}_*\circ[Q']\circ\mathcal{U}.
\]
Since this holds for every HM-pair $(F,Q)$, $A$ is strongly stacked.
\end{proof}

\begin{cor}
Every subfit KC-frame is spatial.
\end{cor}

\begin{proof}
Let $A$ be a subfit KC-frame. Since every KC-frame is $T_1$, the
result follows from \cite[Proposition~2.5.3]{picado2021separation}.
\end{proof}

\begin{prop}
There exists a KC-frame which is strongly stacked and not subfit.
Consequently,
\[
    \mathrm{KC}\not\Rightarrow\mathrm{(sfit)}.
\]
\end{prop}

\begin{proof}
Consider the frame
\[
    A=B\times 2,
\]
introduced above, where $B$ is the complete Boolean algebra considered
in the previous example. Recall that $A$ is $T_1$, strongly stacked,
and non-spatial.

The only open filters of $A$ are
\[
    A
    \qquad\text{and}\qquad
    F_p=\uparrow(0,1)=B\times\{1\},
\]
where $p=(1,0)$ is the unique point of $A$.

For $F=A$, we have
\[
    v_A=w_A=\operatorname{tp}=u_1.
\]
Hence
\[
    [v_A,w_A]=\{u_1\}.
\]

On the other hand, for $F=F_p$, the computations in the previous
example give
\[
    v_{F_p}=w_{F_p}=u_p.
\]
Therefore,
\[
    [v_{F_p},w_{F_p}]=\{u_p\}.
\]

Thus, for every $F\in A^\wedge$, every nucleus in the admissibility
interval $[v_F,w_F]$ is closed. Hence $A$ is a KC-frame.

Finally, $A$ cannot be subfit. Indeed, every $T_1$ subfit frame is
spatial, whereas $A$ is $T_1$ and non-spatial. Consequently, $A$ is
a KC-frame which is strongly stacked and not subfit. In particular,
\[
    \mathrm{KC}\not\Rightarrow\mathrm{(sfit)}.
\]
\end{proof}

\begin{prop}\label{prop:KC}
Every neat strongly stacked frame is KC.
\end{prop}

\begin{proof}
Let $F\in A^\wedge$. Since $A$ is tidy, it is $T_1$. Hence, as $A$ is
strongly stacked, we have
\[
v_F=w_F.
\]
Moreover, of tidiness,
\[
v_F=u_d,
\qquad d=v_F(0).
\]
Therefore,
\[
[v_F,w_F]=\{u_d\}.
\]

Since every nucleus inducing a compact quotient of $A$ belongs to an
admissibility interval $[v_F,w_F]$ for some $F\in A^\wedge$, every such
nucleus is closed. Hence every compact quotient of $A$ is closed, and
therefore $A$ is KC.
\end{proof}

\bibliographystyle{amsalpha}

\bibliography{research2}

\newpage

%\section*{Examples}

%\subsection*{The frame $[0,1]$}
Following the idea presented in \cite{simmons2006regularity} (or, alternatively, in \cite{simmons2004vietoris}), let us consider the linearly ordered frame 
\[
A=[0,1].
\] 
In this case the following observations hold.
\begin{enumerate}
\item For all $a,b\in A$ the implication $(a\succ b)$ is given by
\[
(a\succ b)=\begin{cases}
1, & \text{si } a\leq b,\\
b, & \text{si } a>b.
\end{cases}
\]
\item For $a, x\in A$, the distinguished nuclei $u_a, v_a, w_a$ are given by
\[
u_a(x)=\begin{cases}
  x & \text{si }a\leq x,\\
  a & \text{si }x< a,
\end{cases}, \quad v_a(x)=\begin{cases}
1 & \text{si } x\geq a,\\
a & \text{si } x<a,
\end{cases} \quad w_a(x)=\begin{cases}
1 & \text{si } x>a,\\
a & \text{si } x\leq a.
\end{cases}
\]

\item If $F\subseteq A$ is a filter, then $F$ is always admissible and has the form
\[
F=\nabla(w_a)=(a, 1]\quad \mbox{ o }\quad F=\nabla(v_a)=[a, 1]
\]
for some $a\in A$. In particular, $F\in A^\wedge$ if and only if $F=(a, 1]$.

\item If $F\in A^\wedge$, the adjusted nucleus associated $v_F$ is calculated as
\begin{equation}\label{eq:lm}
v_F(x)=l_m(x)=\begin{cases}
1 & \text{if } x> m,\\
x & \text{if } x\leq m.
\end{cases}
\end{equation}
where $l_m=v_m\wedge w_m$ and $m=\bigwedge F$.
\item For $S=\pt A$ it holds that $S=[0,1)$ and the open sets in $\mathcal{O}S$ are of the form
\[
\mathcal{U}_A(x)=\{p\in S\mid p<x\}=[0,x).
\]
\item The frame $A$ is not $T_1$, since no $p\in S$ is maximal. Consequently, $A$ does not satisfy any stronger separation property in $\Frm$.
\item For $F\in A^\wedge$ the compact saturated set associated $Q\in \mathcal{Q}S$ is of the form $Q=[0,m]$.
\item It holds that $[Q']=v_F$; therefore, $S$ is strongly stacked.
\item Since $[0, x)\in \mathcal{O}S$, then $[x, 1]\in \mathcal{C}S$ and $Q=[0, m]$. Therefore, the compact subsets are not closed; that is, $S$ is not packed.
\item The frame $A$ is not efficient.

\end{enumerate}

The above resume some characteristics about the frame $[0,1]$. Now, we will perform the analysis of the patch constructions of 
$A$ and its point space $S$.

Remember that for $A\in \Frm$ and $S\in \Top$ we have 
\[
\text{Pbase}=\{u_a\wedge v_F\mid a\in A, F\in A^\wedge\}\quad \text{y}\quad \text{pbase}=\{U\cap Q'\mid U\in \mathcal{O}S, Q\in \mathcal{Q}S\}
\]
where Pbase generates the patch frame of $A$ and pbase generates the topology of the patch space of $S$. By \eqref{eq:lm} we have that 
\[
\text{Pbase}=\{u_a\wedge l_m\mid a\in [0,1], m\in [0,1)\}.
\]

For $x\in A$  we have
\[
(u_a\wedge l_m)(x)=u_a(x)\wedge l_m(x).
\]

Let us analyze the value of this map according to the relative position of $x$ and $m$.
\begin{description}
\item[Case 1] If $x<m$, then 
\[
(u_a\wedge l_m)(x)=u_a(x)\wedge x=x.
\]
\item[Case 2] If $x=m$, then 
\[
(u_a\wedge l_m)(x)=u_a(m)\wedge m=x.
\]
\item[Case 3] If $x>m$, then 
\[
(u_a\wedge l_m)(x)=u_a(x)\wedge 1=u_a(x).
\]
\end{description}

In the first two cases, the value of $a$ does not affect the evaluation. Thus, the only non-trivial case occurs when $m<x$, where the value of $(u_a\wedge l_m)(x)$
depends on the relative position of $a$ and $x$. We therefore distinguish the following two subcases.
\begin{description}
\item[Case 3.1] If $m<x<a$, then 
\[
(u_a\wedge l_m)(x)=a.
\]
\item[Case 3.2] If $m<a\leq x$, then 
\[
(u_a\wedge l_m)(x)=x.
\]
\end{description}

From the preceding analysis, we conclude that the non-trivial generators of Pbase correspond precisely to the pairs $(m,a)$ with $m<a$. Consequently, 
the frame $PA$ is determined by the open intervals of the form $(m,a)$, where $m\in\pt A$ and $a\in A$.

We now turn to the question of whether $PA$ satisfies any separation property. To this end, we recall the relationship between the point-free patch construction and its spatial counterpart.

Let $A$ be an arbitrary frame and let $S=\pt A$ be its space of points. The two patch constructions are related by the following commutative diagram:
\[\begin{tikzcd}
	A & PA & NA \\
	{\mathcal{O}S} & {P\mathcal{O}S} & {N\mathcal{O}S} \\
	& {\mathcal{O}^pS} & {\mathcal{O}^fS}
	\arrow[from=1-1, to=1-2]
	\arrow["{\mathcal{U}_A}"', from=1-1, to=2-1]
	\arrow[hook, from=1-2, to=1-3]
	\arrow["{P(\mathcal{U}_A)}", from=1-2, to=2-2]
	\arrow["{N(\mathcal{U}_A)}", from=1-3, to=2-3]
	\arrow[from=2-1, to=2-2]
	\arrow[hook, from=2-1, to=3-2]
	\arrow[hook, from=2-2, to=2-3]
	\arrow["\pi", from=2-2, to=3-2]
	\arrow["\sigma", from=2-3, to=3-3]
	\arrow[hook, from=3-2, to=3-3]
\end{tikzcd}\]
where $\mathcal{O}^fS$ is the Skula topology (or front topology) of $S$. In \cite{sexton2006point} is named
\emph{the patch assembly diagram} and more details about it can also be found in \cite{sexton2003point}.

In our case, $\mathcal{U}_A$ is an isomorphism. Moreover, under $\mathcal{U}_A$, the patch construction turns out to be functorial and, therefore,
$PA\cong P\mathcal{O}S$. Similarly, $NA\cong N\mathcal{O}S$.
To describe explicitly $P\mathcal{O}S$ and $N\mathcal{O}S$, we use the following result, which can be consulted in \cite{AvilaBezhaniMorandiZaldivar2020}.

\begin{prop}\label{Ensamble espacial}
Let $S$ be a partially ordered set.
\begin{enumerate}
  \item If $S$ has no infinite antichains, then $N\mathcal{O}S$ is spatial.
  \item If $S$ is totally ordered, then $N\mathcal{O}S$ is spatial.
\end{enumerate}
\end{prop}

Thus, by Proposition \ref{Ensamble espacial} we have that $N\mathcal{O}S$ is spatial and therefore $N\mathcal{O}S\cong \mathcal{O}^fS$. With all this information,
the patch assembly diagram, in our particular case is as follows.

\[\begin{tikzcd}
	A & PA & NA \\
	{\mathcal{O}S} & {\mathcal{O}^pS} & {\mathcal{O}^fS}
	\arrow[from=1-1, to=1-2]
	\arrow["\cong"', from=1-1, to=2-1]
	\arrow[from=1-2, to=1-3]
	\arrow["\cong", from=1-2, to=2-2]
	\arrow["\cong", from=1-3, to=2-3]
	\arrow[from=2-1, to=2-2]
	\arrow[from=2-2, to=2-3]
\end{tikzcd}\]
Therefore, if we want to study the frame $PA$, we can study the patch topology $\mathcal{O}^pS$. Remember that 
\[
\text{pbase}=\{U\cap Q'\mid U\in \mathcal{O}S, Q\in \mathcal{Q}S\}.
\]
In our case,
\[
\text{pbase}=\{[0, a)\cap [m, 1)\mid a\in [0,1], m\in [0,1)\}=\{(m,a)\}.
\]

Since the topology generated by the intervals $(m,a)$ coincides with the usual topology on $[0,1)$, the space $^pS$ is Hausdorff. Consequently, $\mathcal{O}^pS\cong PA$ satisfies $\mathbf{(H)}$, 
even when $A$ is not $T_1$ nor tidy.

\begin{lem}\label{patchiso}
Let $A$ be a frame, if the morphism $\mathrm{p}U\colon\mathrm{p}A\rightarrow\mathrm{p}\mathcal{O}S$ 
is an isomorphism, then $A$ is strongly stacked.

\end{lem}


\begin{proof}
Consider any open filter $F$ on $A$ and the compact saturated determined by it, say $Q$, 
thus we have the nuclei $\mathrm{v}_{F}$ and $U_{*}[Q']U=Q[U]$, we neee to observe the action of 
$\mathrm{p}U$ on these two nuclei, in order to do that remember that we have a conmutative diagram:



\[\begin{tikzcd}
	A & \mathrm{p}A & NA \\
	{\mathcal{O}S} & {\mathrm{p}
	\mathcal{O}S} & {N\mathcal{O}S} \\
	& {\mathcal{O}^pS} & {\mathcal{O}^fS}
	\arrow[from=1-1, to=1-2]
	\arrow["{\mathcal{U}_A}"', from=1-1, to=2-1]
	\arrow["\iota",hook, from=1-2, to=1-3]
	\arrow["{\mathrm{p}(\mathcal{U}_A)}", from=1-2, to=2-2]
	\arrow["{N(\mathcal{U}_A)}", from=1-3, to=2-3]
	\arrow[from=2-1, to=2-2]
	\arrow[hook, from=2-1, to=3-2]
	\arrow["\imath "',hook, from=2-2, to=2-3]
	\arrow["\pi", from=2-2, to=3-2]
	\arrow["\sigma", from=2-3, to=3-3]
	\arrow[hook, from=3-2, to=3-3]
\end{tikzcd}\]

since $N(\mathcal{U})\iota=\imath\mathrm{p}\mathcal{U}$, that is, $\mathrm{p}\mathcal{U}$ 
is the restriction of $N(\mathcal{U})$ to $\mathrm{p}A$, then 

\[\mathrm{p}\mathcal{U}(\mathrm{v}_{F})=N(\mathcal{U})(\mathrm{v}_{F})=\bigvee\{u_{\mathcal{U}(v_{F})}\wedge v_{\mathcal{U}(x)}\mid x\in A \}\]

since the morphism $\mathcal{U}$ transforms open filters we have 



\[\bigvee\{u_{\mathcal{U}(v_{F})}\wedge v_{\mathcal{U}(x)}\mid x\in A \}=\bigvee\{u_{v_{\mathcal{U}F}}\wedge v_{\mathcal{U}(x)}\mid x\in A \}\]


But this representation is the corresponding to the nucleus $\mathrm{v}_{\mathcal{U}F}$.
\todo{Completar esta prueba}

\todo{Considerar incluir información sobre los parches}

\end{proof}

\end{document}
